\chapter[India: Traditions Beyond One Box]{Why India's Religious Traditions Resist a Single Box}
\section{A Picture on the Wall}
First, pick up an object: a picture.

It is about the size of a magazine, printed on glossy paper and mounted in a cheap gold-colored plastic frame. It shows a plump god with an elephant's head and a human body, sitting on a lotus. His four arms hold an axe, a rope, and a bowl of sweets, while one hand is raised toward you, palm outward, as if saying, "Don't be afraid." A mouse crouches at his feet. A small line in the lower-right corner names a printing works in Mumbai.

Pictures like this are everywhere in India. They hang on the walls of corner grocery shops, above taxi dashboards, and beside the tills in small restaurants. The figure is the elephant-headed god Ganesha, said to "remove obstacles." When opening a shop, taking an examination, setting off on a long journey, or signing a contract, many people think of him first.

But first look at where this picture hangs. It is not the only image on the grocery shop's wall. To its left is a photograph of a white-bearded old man sitting cross-legged: the guru whom the shopkeeper's grandfather followed during his lifetime. To its right hangs a calendar showing the domes of the Sikh Golden Temple, because the company that printed it is in Punjab. A sticker of a Bollywood actress adorns the cash register---the shopkeeper's son insisted on putting it there.

Now a question that will follow us throughout this chapter: to "which religion" does that picture of Ganesha belong?

You might answer without thinking: Hinduism, of course. That is what textbooks, encyclopedias, and the news call it. But notice something subtle: the shopkeeper may never in his life have used the word "Hinduism" for what he does. Every morning he offers Ganesha fresh flowers and clean water; when trouble comes, he visits a goddess's temple outside the village; a priest must consult the stars to choose his child's wedding date; the anniversary of his father's death requires another kind of ceremony. The rules, deities, and languages of recitation differ from one practice to another. They resemble layers placed in a large basket at different moments in history. The basket called "Hinduism" was itself made only later.

In this chapter we will unpack that basket. Why is this ancient, enormous, still-living tradition of the South Asian subcontinent so difficult to fit into a box labeled "religion"? Who made the box? Does putting the tradition inside help us see it clearly, or accidentally squash it out of shape?

\section{The Ganges Steps Before Dawn}
Now come with me to a particular scene.

Time: five in the morning. Place: Varanasi, a city in northern India, on the stone steps beside the Ganges. Locally these are called ghats: rows of stone steps descending from the bank into the water. In Indian tradition, the river is herself a goddess.

Day has not yet broken. The air is cool and damp, carrying a mixture of sandalwood, jasmine, river mud, and a little smoke. People are already on the steps.

A group of elderly women wrapped in saris descend step by step to the water, press their palms together, and immerse themselves in the waist-deep river. The water is cold, but they keep murmuring prayers. On a wooden platform above them, a priest raises a many-tiered brass lamp and slowly traces circles toward the river. Bells, chanting, and drums sound together. This is a ceremony welcoming the sunrise; it was performed last night and will be performed again tomorrow night. The same movements, people say, have been repeated for generation after generation.

Walk a few hundred meters to the left and the scene changes completely. Heavy smoke hangs over another flight of steps, where several piles of wood are burning. This is the riverside cremation ground; the bodies being burned belong to people who died yesterday. Many families bring relatives here from far away, believing that ending a life beside this river will spare the soul much of the suffering of rebirth. Family members sit silently beside the pyres. The wood sellers charge by weight.

Farther to the right, at the entrance to an alley above the steps, a young man wearing a head wrap spreads out a mat and performs the Sun Salutation. Its name is almost the same as that of the priest's ceremony a few hundred meters away: both "greet the sun." Yet he may have no interest in how many circles the priest traces with his brass lamp. A chai vendor emerges from the alley, tea bubbling in his tin kettle. Above his head, a string of mango leaves and a little Ganesha plaque hang from a doorway. A visitor from Europe sits in the bow of a boat, recording everything on a phone, their expression a mixture of awe and puzzlement.

Stand here for a moment and count: goddess, brass lamp, cremation, yoga, chai, divine images, visitor. They crowd into a few hundred meters of riverbank, happening simultaneously without contradiction. Nobody directs everyone else, and no "instruction manual" specifies how these things should fit together.

This scene gives flesh to the question of our chapter.

\section{One Question: Is This "One Religion"?}
Let us first set out the tension, without rushing to an answer.

Suppose a census taker now walks onto the steps, carrying a form and asking everyone in turn: "What is your religion?" The form has just one field, and each person can enter only one word.

What should an elderly woman write? She worships the Ganges goddess, the Ganesha in her home, and a village goddess; her grandson has married a woman from a neighboring village who follows a different god. What should the descendant of a cremation worker write? His family has done this work for generations because birth assigned them to society's lowest rank. Yet he is the one who takes others through the final stage toward "liberation": the holiest ceremony is performed by the person considered "lowest." What should the young yoga practitioner write? He might say he is not superstitious, merely exercising. What should the priest write? The scriptures he memorizes are three thousand years old, yet they say nothing about the god his great-grandmother worshipped.

Still the form demands: one box, one word.

There are several layers of difficulty here, each deeper than the last.

The first: this tradition has no founder. Buddhism has Shakyamuni, Christianity has Jesus, and Islam has Muhammad. But for this collection of things beside the Ganges, there is no name to enter under "founder." It grew layer by layer over thousands of years, like an old city that keeps adding buildings without ever demolishing its older houses.

The second: it has no single, universally accepted "Bible." It has the Vedas, but after the Vedas came the Upanishads, the two great epics, the Puranas, and poetry in many regional languages. These accumulated in layers, not entirely consistent with one another, without anyone announcing which book must be the final authority.

The third, and most fundamental: the "religion" box we are holding has a problematic shape. It was broadly modeled on certain West Asian traditions: a founder, scriptures, an initiation ceremony, membership in one excluding membership in another, and a core list of propositions stating "what you believe." Measuring the traditions beside the Ganges with this box is like trying to put a forest in a wardrobe. Some things fit; many more stick out.

The real question now surfaces: \textbf{When we say "Hinduism," are we describing something that already exists, or forcing a forest into a box brought from elsewhere?}

Before reading on, take a position of your own: without the umbrella name "Hinduism," would everything beside the Ganges still count as "one thing"?

\section{Six or Seven Answers Beside One River}
Whether this counts as one religion depends on whom you ask. Let us hear the voices beside the river one by one---at least six of them---before allowing ourselves a judgment.

\textbf{The priest's voice: this is a river that has flowed for three thousand years.}

The priest will tell you that everything has a source. The scriptures he recites are called the Vedas. Their oldest parts began to be compiled more than three thousand years ago and have been passed down orally, memorized word by word and sound by sound, to the present day. Before recording machines existed, this was the most precise technology of preservation. To him, everything that followed---images, temples, yoga, village goddesses---is a tributary of this great river. Its waters may divide into countless channels, but they come from the same stream. This position has a particular name. Many followers call their tradition the "eternal law," meaning that nobody invented it in a particular year: it has always been there.

\textbf{The grandmother's voice: I only know whom I worship.}

Visit a southern Indian village and ask an elderly grandmother draping red gauze over a goddess's image, "What is your religion?" She probably will not know how to answer. She worships her village's goddess---perhaps Mariamman, associated with rain and disease and deeply revered in Tamil villages. This goddess has her own stories, feast days, and favorite offerings, with no direct connection to the Vedas the priest recites. For the grandmother, religion is not a list of beliefs but the actions of a day: drawing auspicious patterns in rice flour outside the door at dawn, lighting an oil lamp before the shrine at dusk, and consulting the goddess about family matters large and small. This is not a "budget version" of religion. Quite the opposite: this is what the religious lives of the vast majority of people worldwide look like.

\textbf{The philosopher's voice: a divine image is only a finger pointing at the moon.}

About twelve hundred years ago, a thinker named Shankara traveled across India debating others. His argument, in broad terms, was this: the countless deities and all the distinctions in the world are ultimately appearances. Beneath them lies a single ultimate reality, which he called Brahman, and the "self" deep within each person is fundamentally not separate from it. Worshipping an image is not wrong, but one must understand that an image is like a finger pointing at the moon: stare at the finger and you miss the moon. This thought derives from the Upanishads, a collection of philosophical dialogues that took shape after the Vedas, beginning roughly twenty-five hundred years ago. In other words, while Greek philosophers were debating the ultimate origin of the world, people in the forests beside the Ganges were debating questions of comparable depth.

\textbf{The weaver's voice: God is not in the stone.}

In fifteenth-century Varanasi lived Kabir, a poet who came from a weaving background. Born into a Muslim family, he nevertheless sang in the language of Indian traditions. The gist of his poetry was this: you rush to temples and bow, but God is not in the stone; you say God is in the mosque, but God is not there either; God is in the sincere heart. He mocked priestly rituals and empty dogma, and openly opposed ranking people by birth. This movement of "loving God directly, without relying on priests" is called \textbf{bhakti}, or devotion. Extending from south to north over a thousand years, it produced many poets of humble birth, including women. Ironically, Kabir's poems were eventually included even in Sikh scripture: a man who rejected all "sectarian boxes" ended up being claimed by several of them.

\textbf{The skeptic's voice: there is no afterlife at all.}

We must insert a counterexample that is easily overlooked. Many people assume that in India everyone believes in gods and rebirth. This is wrong. Radical skeptics have always stood within Indian traditions. An ancient school called Charvaka, or Lokayata, openly maintained that there are no gods, no soul, and no afterlife, and that sacrifice is a business through which priests trick people into feeding them. These views were recorded by its opponents more than two thousand years ago. Among the six major schools of Indian philosophy, Nyaya specializes in logic and the rules of argument, with a sophistication equal to that of logic anywhere; Samkhya scarcely needs a god to explain the world. In other words, the impression "India equals mysticism equals piety" is a coat of paint applied by outsiders. Skepticism has a history on this land as long as that of devotion.

\textbf{The "untouchable's" voice: there is no place for me in this extended family.}

In the twentieth century, B. R. Ambedkar stepped forward with the most uncomfortable challenge. Born into an "outcaste" family trampled into traditional society's lowest rank, he was forbidden by the old rules even to touch his school's well as a child. Through extraordinary perseverance he studied abroad, became a legal expert, and later led the drafting of India's Constitution. Throughout his life he asked: if "Hinduism" really is one extended family, why are some of its members declared people whom others may not touch, eat with, or admit to the same temple? His public declaration amounted to this: I could not choose my birth, but I will not die within this system as an "untouchable." In 1956 he led hundreds of thousands of followers in a collective conversion to Buddhism, leaving the box through action. We must listen honestly to this voice. It reminds us that the attractive word "diversity" can sometimes sound, to those at the bottom, like a sentence declaring: "The bottom is where you belong."

\textbf{The census official's voice: without classification, government is impossible.}

Finally, let us state the form-carrying official's case in full. This is our chapter's "steelmanning" exercise: deliberately adopting a position you disagree with and making its argument as strong as possible.

The census official did not appear from nowhere. In the late nineteenth century, the British colonial government ruling India began nationwide censuses requiring everyone to report a single "religion." Put the case as fully as possible: government needs numbers, and numbers need categories. Should schools provide classes for different groups? Which rules should settle legal disputes? Military recruitment, taxation, electoral boundaries---which of these can do without sorting people into categories? Nor is this a peculiar habit invented by colonizers. Every country's statistical offices do the same today. Requiring "one person, one box" is not malice but a practical necessity of modern administration. This is one real origin of the "Hinduism" box: not a desire to squash the forest, but the need for boxes on a form.

We have heard them now: seven voices, seven accounts of "what it is." Notice that their disagreements are not entirely about facts. The priest and the grandmother see the same river, but what they mean by "tradition" belongs to quite different levels.

\subsection{Four Key Words, Four Sets of Emphases}
There is another, less visible reason these voices clash: they use the same words but place the emphasis differently. South Asian traditions have four key terms, each with multiple meanings.

\textbf{Dharma}, also rendered in Chinese by a phonetic transliteration. It can mean the order of the universe, by which the sun rises and sets; the duties of your social position, with sons and kings each having their dharma; justice; or simply "religion" itself. A Hindu who speaks of "keeping one's dharma" and a Buddhist who speaks of "hearing the Buddha's Dharma" use the same word with two meanings. When people argue over whose dharma is the true dharma, they may not even be speaking within the same sentence to begin with.

\textbf{Karma}. Its literal meaning is "action," extended to the consequences actions produce. But its uses differ enormously. For some, karma resembles an exact moral ledger: suffering in this life repays debts from a previous one. This can make suffering seem "reasonable," which is precisely its danger, because it may imply that poor people deserve their poverty. Buddhism redefined karma: what determines it is not the act itself but the intention behind it. The same cut made by a surgeon and by a murderer produces entirely different karma. One word, two systems of moral physics.

\textbf{Samsara}, the cycle of rebirth. Many traditions share the idea that life passes from one existence to another, but they have argued for thousands of years about what passes between them. Most Indian traditions posit an unchanging "self" undergoing rebirth. Buddhism, by contrast, makes the absence of a fixed, unchanging self foundational. Rebirth continues, but what continues is not a substantial soul. This is one of the most finely drawn disagreements in religious history.

\textbf{Liberation} (moksha/nirvana). If rebirth is a cycle, liberation means leaving it. Different traditions point to different exits. Shankara says liberation is recognizing that the self and Brahman are one. Bhakti poets say it is remaining forever near the god one loves, even as a gatekeeper. Buddhist nirvana means extinguishing the fires of greed, hatred, and delusion, relinquishing even the thought "I have reached the farther shore." All speak of "salvation," but what is saved differs.

These four terms resemble four characters with multiple pronunciations. Once you understand their many meanings, you can see that even the most central vocabulary of "Indian religious tradition" is plural.

\section[Thinking Bridge: Family Resemblance]{Thinking Bridge: Seek "Family Resemblance," Not an Essence}
To approach this many-layered collection of traditions, we have a portable thinking tool from the twentieth-century philosopher Ludwig Wittgenstein.

He posed a seemingly ordinary question: what exactly does the word "game" mean? Chess is a game, football is a game, cards are a game, and children's hide-and-seek is a game. Try to find the feature \textbf{shared by all games and possessed only by games}. Winning and losing? What about a child tossing a ball alone? Rules? Children invent rules as they chase one another. Entertainment? Playing chess is work for a professional. You discover that the essence cannot be found. Wittgenstein said games are connected through \textbf{family resemblance}. In a large family, the second child has the father's eyes, the third has the mother's chin, and the first and fourth have similar personalities. Yet there is no feature that every family member possesses and every outsider lacks. What makes them a family is not a "family essence" but a web of overlapping similarities.

Now bring this tool to the steps beside the Ganges:

\begin{itemize}
\item The Vedic priest's fire sacrifice and flowers offered to an image in a temple share the act of "offering";
\item Temple flowers and the grandmother's rice-flour patterns share the intention of "welcoming a deity with beauty";
\item Shankara's philosophy and the Upanishadic dialogues share the concepts of Brahman and self;
\item The Upanishads and bhakti poetry share an ambivalent relationship with the Vedas.
\end{itemize}
Each neighboring pair of traditions overlaps, yet you cannot find one thing shared by fire sacrifice, village deities, Shankara, Kabir, and the young person on a yoga mat. It is just like "games." So the question "What is Hinduism?" may itself be wrongly framed: it presumes an essence waiting to be discovered. Ask instead, "What overlapping similarities connect these traditions?" The answer immediately becomes richer and more honest.

You can take this tool with you anywhere. "What is Chinese culture?" "What is art?" "What is a true friend?" These are all questions of the same kind. A concept whose examples yield no common essence, however thoroughly you search, is likely a family-resemblance concept. When you meet one, stop asking about its essence and ask how its members resemble one another.

But remember the tool's limits. Family resemblance tells you \textbf{what a group looks like}, not \textbf{where its boundaries lie} or \textbf{who is shut outside them}. When people like Ambedkar are told, "You too belong to this extended family," the warm-sounding phrase may actually demand that they accept the lowest seat assigned to them. Resemblance explains shape, not power. We will return to this problem.

\section{"I Don't Know" Hidden in the Oldest Scriptures}
Now let us read a short primary-source passage. To understand why this tradition resists the box, the best approach is not to listen to outsiders describe it but to open its oldest texts and see what they say.

The Rig Veda is the oldest surviving Indian sacred text, compiled more than three thousand years ago in Sanskrit and still part of humanity's shared public-domain heritage. Its tenth book contains a famous hymn often called the "Hymn of Nonbeing" (Rig Veda 10.129). It sings of creation, but in an unexpected way. Its general meaning is:

\begin{quote}
Then there was neither "being" nor "nonbeing," neither the sky nor what lay above it. What covered everything? Where? Under whose protection? ... There was no death, nor immortality, no distinction between day and night. ... Who truly knows? Who can explain it here? From where did all this arise? From where did this creation come? The gods appeared only after creation---so who knows where it really came from?
\end{quote}
Read the last two sentences again. A creation hymn in a sacred scripture ends, astonishingly, with this: perhaps nobody knows; even the gods arrived too late to witness the beginning.

This is a highly unusual stance for a religious scripture. We are accustomed to scriptures opening with firm answers. Here, instead, is a sacred hymn that writes doubt into its own core. It does not forbid questions; it makes questioning to the very end part of the sacred text itself.

The same Rig Veda contains another line repeatedly quoted by later generations (1.164). Its commonly conveyed sense is:

\begin{quote}
What is real is one; the wise call it by many names.
\end{quote}
One or many? This ancient verse answers both sides at once: the divine names may number in their thousands, while reality is one. The line later became one of this tradition's most frequent explanations of its own diversity. It suggests that resisting a single box was not a fault that developed later: diversity was written on the cover from the earliest texts onward.

This material also corrects a common misreading. Many people assume that internal diversity arose because later generations muddled something originally pure. The evidence points the other way: "perhaps nobody knows" was already written on the oldest page. Doubt is not foreign contamination; it came with the original package.

\section{One Word, Three Histories}
We can now formally compare explanations. They all address the same fact: why does "Hinduism" today encompass such radically different things as fire sacrifices, village deities, philosophical debates, devotional singing, and young people on yoga mats who believe in no god? There are at least three genuinely distinct, nonequivalent explanations.

\textbf{Explanation one: outsiders made the box.}

This explanation points out that "Hindu" was originally a geographical rather than religious name. Persians used a term derived from Sindhu, the ancient name of the Indus, for the people and land beyond the river. The word subsequently traveled through various languages, retaining the meaning "on the other side of the Indus." Turning "the people of India" into "Hinduism," a word with an "-ism" suffix, happened much later, roughly in eighteenth- and nineteenth-century English writing. Late-nineteenth-century colonial censuses then made membership in exactly one religion an administrative requirement. Its strength: it explains why the word is so young and why the box happens to resemble a field on an administrative form. Its limitation: it cannot explain the real connections inside the box. Why were Kabir's poems included in Sikh scripture? Why do temples separated by great distances from north to south worship the same Shiva? If everything is merely an outsider's label, where do these internal cross-references come from?

\textbf{Explanation two: it was always a continuous civilization.}

This explanation responds directly: diversity does not mean an arbitrary assemblage. From the Vedas to the Upanishads, from the epics Mahabharata and Ramayana to regional bhakti poetry, texts quote one another and enter into layered conversations. Pilgrimage networks spanning the subcontinent connect northern mountain shrines and southern seaside temples on one map. More importantly, believers themselves have always felt part of "the same thing," calling it the "eternal law." Its strength: it respects participants' self-understanding and explains the real cultural and pilgrimage networks. Its limitation: presenting tradition as a single whole can easily smooth away internal disputes. Charvaka skepticism, Kabir's mockery of priests, and Ambedkar's departure are all turned down in volume within the story of a harmonious extended family. Who benefits most from saying, "We have always been one family"? Often precisely those at the top of the household.

\textbf{Explanation three: it is the product of centuries of layering and interpenetration.}

The third explanation changes the timescale. It says: do not ask what the tradition originally was; ask how it grew. Historically, elite traditions written in Sanskrit interacted over long periods with regional traditions. A local goddess might be welcomed into Sanskrit mythology and "promoted" to an incarnation of a major deity; conversely, a village might reshape a festival from the wider tradition into something local. Sociologists call the movement of local traditions toward the larger tradition "Sanskritization." Its strength: it closely fits what fieldwork reveals. In some villages, for example, Mariamman's story now marries her to a deity within Shiva's divine world, leaving clear signs of grafting. Its limitation: explaining everything risks explaining nothing. If everything can be called "layered interaction," it becomes difficult to specify when a particular layer was added, by whom, and in whose interests.

Each explanation holds part of the truth. What matters is not choosing one but taking away a method: when an explanation claims to account for everything, it has probably discarded something in advance.

\subsection{Caste on Paper and Caste on the Ground}
Comparing explanations also brings an unavoidable test: caste. It is another especially difficult case for boxes, because the system on paper never quite matches life on the ground.

The paper version is orderly. A famous passage in the "Hymn of the Cosmic Person" in the Rig Veda's tenth book (10.90) says, in broad terms, that all things emerged from the body of the primordial giant: brahmins, or priests, from his mouth; royal warriors from his arms; common people from his thighs; and the lowest-ranking servants from his feet. Later legal compilations, such as the Manusmriti and related "dharma-sutra" texts, turned these four ranks into detailed rules governing who could read sacred texts, who could marry whom, and who could pursue which occupation. Read only these texts and you might imagine Indian society as a machine operating according to the same blueprint for two thousand years.

The ground-level version is quite different. Everyday life is governed not by those four large categories, or varnas, but by thousands of particular birth groups, or jatis. Often tied to locality and occupation, their rankings vary completely across regions. A group ranked highly in one state might be unknown in its neighbor. Historical records also contain many examples of groups collectively changing status, changing occupations, and rewriting genealogies to move upward. Dharma-sutra texts describe what ought to happen, not what actually happens. Mistaking prescription for reality is among the easiest mistakes to make when reading Indian sources.

More importantly, criticism of caste is as old as the system itself. The Dhammapada is one of the earliest Buddhist verse collections, a Pali scripture in the public domain. Its verse 396 states, in unambiguous general terms:

\begin{quote}
Birth does not make a person an "outcaste," nor does birth make a brahmin; actions make an "outcaste," and actions make a brahmin.
\end{quote}
This is a public rebuttal from more than two thousand years ago: status rests on conduct, not bloodline. In other words, even the tradition's oldest internal texts reject the claim that caste is an unquestionable part of Indian tradition. People within the tradition challenged it from the beginning. This is the same pattern we have seen throughout: every "box" in this tradition contains people dismantling it.

\section[Religion Beyond a Creed]{Further Challenge: Religion Is More Than a Creed}
We now introduce the chapter's weightiest theoretical tool, from a discipline you may not yet have encountered: religious studies---not theology. Theology speaks from within a particular tradition; religious studies stands outside and comparatively examines the world's religions.

The twentieth-century scholar Ninian Smart proposed a simple but powerful idea. We tend to imagine "religion" as a list of beliefs---whether you believe in God or an afterlife---as if its core were a set of propositions to check off. Yet observing actual religious life reveals at least \textbf{seven dimensions} operating together:

\begin{itemize}
\item \textbf{Ritual}: what people do (fire sacrifices, flower offerings, river bathing, lighting lamps);
\item \textbf{Narrative and myth}: the stories people tell (how Ganesha acquired his elephant's head);
\item \textbf{Experience and emotion}: what people feel (the shiver when Ganges water touches the skin);
\item \textbf{Doctrine and philosophy}: what people argue (reasoning about the nonduality of Brahman and self);
\item \textbf{Ethics and law}: what people prescribe (what your dharma is);
\item \textbf{Society and organization}: who occupies which position (priests, temples, castes, religious orders);
\item \textbf{Material things}: what objects people use (divine images, brass lamps, rice-flour designs, that picture).
\end{itemize}
Look back at the Ganges through this framework and things become clear. People there have not failed to complete a belief checklist; their religious lives were never centered on one. The grandmother's religion is full in its ritual, material, and emotional dimensions, nearly empty in its doctrinal one. Shankara's religion lies almost entirely in doctrine. Kabir's centers on experience and deliberately undermines the social dimension. Testing them only on "what you believe" is like judging a painting solely by its volume.

Smart's framework also reveals something important: the box equating religion with a creed has its own history. After the European Reformation, in an age when monotheistic traditions debated one another and wrote statements of faith, that shape gradually became the standard for religion, then spread worldwide through colonialism, education, and forms. In other words, \textbf{Indian traditions are not oddly shaped; the ruler measuring them is biased}. Apply that ruler elsewhere and Japanese Shinto---where many people follow Shinto rites at birth, marry in a church, have Buddhist funerals, and describe themselves as nonreligious---and Chinese folk practices of burning incense and honoring ancestors all become "defective religions." A crooked ruler makes the whole world appear crooked.

Of course, theory has limits too. Seven dimensions can lay religious life out before us, but they cannot establish which matters most to a participant. For the grandmother, the rice-flour design outside her door may outweigh every doctrine she has ever heard put together. Completing the map is not the same as completing the journey.

\section[Today: Forms, Yoga, and Politics]{Echoes Today: Forms, Yoga Mats, and Political Slogans}
This ancient question is sounding around us right now.

First, India itself. The Constitution that took effect in 1950 expressly abolished "untouchability" in Article 17, making discrimination against people of "outcaste" birth a crime under law. From the same period, India also reserved a proportion of places in education and public employment for historically oppressed groups. Whether to retain those reservations, how large they should be, and how long they should last remain among India's most fiercely debated public issues. That is the factual progress. There is also a real contemporary controversy: in recent decades, a powerful political current has sought to turn "Hinduism" from a loose umbrella term into a political identity with firm boundaries and a unified position---welding this chapter's box shut and carrying it into the streets. Supporters call this recognition of the majority's tradition; critics say it will squeeze both internal diversity and groups outside the tradition, including Muslims, out of public life. The verified fact is that this debate exists and has enormous influence; where it will ultimately take the tradition remains an open question. Its existence demonstrates that our "battle over the box" is no antiquarian exercise confined to a study.

Now something closer to home. Yoga studios are everywhere in your city. Yoga was originally a practice of one philosophical school discussed in this chapter, the Yoga school, with a complete worldview and purpose. As it spread globally, the worldview was removed and the movements remained, becoming a form of exercise. Is that good or bad? Supporters say the essence naturally travels, so why bind it to everything else? Critics say dismantling a tradition into saleable components is itself a new form of boxing it up. The same story occurs with mindfulness and Zen meditation. You need not take sides immediately, but you can now see the process: some people make boxes, some dismantle them, and some sell them in pieces.

Finally, return to our own lives. At Qingming, the Tomb-Sweeping Festival, you accompany your family to ancestral graves, burn paper offerings, bow, and set out food. If someone asks your religion, you will probably answer, "I am not religious." Yet you have just completed an entire series of rituals, supported by narratives about where ancestors are, emotions of remembrance, an ethic of filial respect, and the social organization of a family. In Smart's framework, your religious life is not empty; it simply has no suitable field on the form. The debate over whether Chinese people have religious beliefs is the same kind of question as whether Hinduism is one religion: the box comes from elsewhere, while the life is your own. Understanding the grandmother beside the Ganges also helps you understand yourself at Qingming.

\section[Your Question: Who Decides the Box?]{Your Question: Who Decides Whether We Need the Box?}
Return to the grocery shop and the picture of Ganesha on its wall.

We now know that the tradition behind it has no founder and no single scripture, and that its oldest hymn ends with "perhaps nobody knows." It contains the devout, philosophers, and thoroughgoing skeptics; a system ranking human beings, and resistance as old as that system. The name "Hinduism" is partly a geographical term and partly a product of modern forms, yet hundreds of millions use it to name themselves.

Here, then, is the chapter's final question: \textbf{Does a tradition like this need a name---a box---at all?}

Give your answer and your reasons. Before deciding, weigh these considerations once more:

Boxes have their uses. You have heard the census official's argument: without classification there are no numbers, and without numbers no law, education, or rights. To protect people pushed to the bottom, India's Constitution first had to identify them as categories. A name is also recognition: a tradition without a name can often be effectively nonexistent in the wider world.

Boxes have their costs too. Every boundary is both an admission pass and an expulsion order. Whoever determines the box's shape holds the power to say, "You do not count." You have heard both Ambedkar's reasons for leaving the box and others' reasons for wanting to weld it shut.

So does putting a forest into a wardrobe protect it or prune it? At the moment a tradition is named, classified, and entered on a form, has it been understood, or has its misunderstanding only just begun? If you designed the form, would you retain the "religion" field? If so, how many words would you allow a person to enter?

There is no standard answer. But remember: the next time a form or a conversation demands that you explain "what you are" in one word, you already know that an entire Ganges lies hidden behind that box.
